%% abtex2-modelo-relatorio-tecnico.tex, v-1.9.6 laurocesar
%% Copyright 2012-2016 by abnTeX2 group at http://www.abntex.net.br/ 
%%
%% This work may be distributed and/or modified under the
%% conditions of the LaTeX Project Public License, either version 1.3
%% of this license or (at your option) any later version.
%% The latest version of this license is in
%%   http://www.latex-project.org/lppl.txt
%% and version 1.3 or later is part of all distributions of LaTeX
%% version 2005/12/01 or later.
%%
%% This work has the LPPL maintenance status `maintained'.
%% 
%% The Current Maintainer of this work is the abnTeX2 team, led
%% by Lauro César Araujo. Further information are available on 
%% http://www.abntex.net.br/
%%
%% This work consists of the files abntex2-modelo-relatorio-tecnico.tex,
%% abntex2-modelo-include-comandos and abntex2-modelo-references.bib
%%

% ------------------------------------------------------------------------
% ------------------------------------------------------------------------
% abnTeX2: Modelo de Relatório Técnico/Acadêmico em conformidade com 
% ABNT NBR 10719:2015 Informação e documentação - Relatório técnico e/ou
% científico - Apresentação
% ------------------------------------------------------------------------ 
% ------------------------------------------------------------------------

\documentclass[
	% -- opções da classe memoir --
	12pt,				% tamanho da fonte
	openright,			% capítulos começam em pág ímpar (insere página vazia caso preciso)
	%twoside,			% para impressão em recto e verso. Oposto a oneside
	a4paper,			% tamanho do papel. 
	% -- opções da classe abntex2 --
	%chapter=TITLE,		% títulos de capítulos convertidos em letras maiúsculas
	%section=TITLE,		% títulos de seções convertidos em letras maiúsculas
	%subsection=TITLE,	% títulos de subseções convertidos em letras maiúsculas
	%subsubsection=TITLE,% títulos de subsubseções convertidos em letras maiúsculas
	% -- opções do pacote babel --
	english,			% idioma adicional para hifenização
	french,				% idioma adicional para hifenização
	spanish,			% idioma adicional para hifenização
	brazil,				% o último idioma é o principal do documento
	]{abntex2}



% PACOTES

% ---
% Pacotes fundamentais 
\usepackage{lmodern}			% Usa a fonte Latin Modern
\usepackage[T1]{fontenc}		% Selecao de codigos de fonte.
\usepackage[utf8]{inputenc}		% Codificacao do documento (conversão automática dos acentos)
\usepackage{indentfirst}		% Indenta o primeiro parágrafo de cada seção.
\usepackage{color}				% Controle das cores
\usepackage{graphicx}			% Inclusão de gráficos
\usepackage{microtype} 			% para melhorias de justificação
% ---
\usepackage{tabularx}
% Pacotes adicionais, usados no anexo do modelo de folha de identificação
\usepackage{multicol}
\usepackage{multirow}
\usepackage{vcell}
\usepackage{float} % para utilizar o H de tabelas e figuras e mante-las no lugar
% ---

%----------------- PARA INSERIR ARQUIVOS DE CODIGOS -----------------%
\usepackage{xcolor}
% Definindo novas cores
\definecolor{verde}{rgb}{0,0.5,0}
% Configurando layout para mostrar codigos C++
\usepackage{listings}
\lstset{
  language=Verilog, %indicar a linguagem utilizada
  basicstyle=\ttfamily\tiny, 
  keywordstyle=\color{blue}, 
  stringstyle=\color{verde}, 
  commentstyle=\color{gray}, 
  extendedchars=true, 
  showspaces=false, 
  showstringspaces=false, 
  numbers=left,
  numberstyle=\tiny,
  breaklines=true, 
  backgroundcolor=\color{green!10},
  breakautoindent=true, 
  captionpos=b,
  xleftmargin=0pt,
}
% -----------------% FIM INSERIR ARQUIVOs DE CODIGO -----------------%

% Pacotes de citações
\usepackage[brazilian,hyperpageref]{backref}	 % Paginas com as citações na bibl
\usepackage[num]{abntex2cite}	% Citações padrão ABNT
\usepackage[utf8]{inputenc}
\usepackage{tcolorbox}
\usepackage{amsmath}
\usepackage{amsfonts}
\usepackage{amssymb}
\citebrackets() %citação numérica entre colchetes
% --- 


% CONFIGURAÇÕES DE PACOTES

% Configurações do pacote backref
% Usado sem a opção hyperpageref de backref
\renewcommand{\backrefpagesname}{Citado na(s) página(s):~}
% Texto padrão antes do número das páginas
\renewcommand{\backref}{}
% Define os textos da citação
\renewcommand*{\backrefalt}[4]{
	\ifcase #1 %
		Nenhuma citação no texto.%
	\or
		Citado na página #2.%
	\else
		Citado #1 vezes nas páginas #2.%
	\fi}%
% ---


% Informações de dados para CAPA e FOLHA DE ROSTO
\titulo{Trabalho computacional 1 - Cálculo Numérico}
\autor{
    Gabriel Augusto Mendonça Beloso RA: 148179 \and
    João Vitor Mota do Amaral RA: 170572 \and
    João Marcelo Constantino Nordemann RA: 168789 \and
    Rafael dos Santos Parisi RA: 148418
} %NAO PREENCHER
\local{São José dos Campos - Brasil}
\data{Setembro de 2026}
\instituicao{
  Docente: Prof. Dr. Luiz Leduíno de Salles Neto
  \par
  Universidade Federal de São Paulo - UNIFESP
  \par
  Instituto de Ciência e Tecnologia - Campus São José dos Campos
}
\tipotrabalho{Relatório técnico}
\preambulo{Relatório apresentado à Universidade Federal de São Paulo como parte dos requisitos para aprovação na disciplina de Calculo Numérico.}
% ---

% Configurações de aparência do PDF final

% alterando o aspecto da cor azul
\definecolor{blue}{RGB}{41,5,195}

% informações do PDF
\makeatletter
\hypersetup{
     	%pagebackref=true,
		pdftitle={\@title}, 
		pdfauthor={\@author},
    	pdfsubject={\imprimirpreambulo},
	    pdfcreator={LaTeX with abnTeX2},
		pdfkeywords={abnt}{latex}{abntex}{abntex2}{relatório técnico}, 
		colorlinks=true,       		% false: boxed links; true: colored links
    	linkcolor=blue,          	% color of internal links
    	citecolor=blue,        		% color of links to bibliography
    	filecolor=magenta,      		% color of file links
		urlcolor=blue,
		bookmarksdepth=4
}
\makeatother
% --- 


% Espaçamentos entre linhas e parágrafos 

% O tamanho do parágrafo é dado por:
\setlength{\parindent}{1.3cm}

% Controle do espaçamento entre um parágrafo e outro:
\setlength{\parskip}{0.2cm}  % tente também \onelineskip

% compila o indice
\makeindex
% ---



% ------------------------------------------------
% Início do documento
\begin{document}

% Seleciona o idioma do documento (conforme pacotes do babel)
%\selectlanguage{english}
\selectlanguage{brazil}

% Retira espaço extra obsoleto entre as frases.
\frenchspacing 

% ----------------------------------------------------------
% ELEMENTOS PRÉ-TEXTUAIS
% ----------------------------------------------------------
% \pretextual

% ---
% Capa
\imprimircapa
% ---

% ---
% Folha de rosto
\imprimirfolhaderosto*
% ---

%------LISTAS SÃO OPCIONAIS---------
% inserir lista de ilustrações
\pdfbookmark[0]{\listfigurename}{lof}
\listoffigures*
\clearpage

% inserir lista de tabelas
\pdfbookmark[0]{\listtablename}{lot}
\listoftables*
\clearpage
% ---

%SUMÁRIO É OBRIGATÓRIO
% inserir o sumario
\pdfbookmark[0]{\contentsname}{toc}
\tableofcontents*
% ---


% ----------------------------------------------------------
% ELEMENTOS TEXTUAIS
\textual

% ----------------------------------------------------------

%------------------- MODELO DE FIGURA -----------------------------
%-Modelo da \autoref{fig:esquematico}:
%-\begin{figure}[H]
%-\centering 
%-\caption{Esquematico} \label{fig:esquematico}
%-\includegraphics[scale=0.3]{esquematico.png}
%-\legend{Fonte: Computer Organization and Design \cite{patterson2007} }
%-\end{figure}


%------------------- MODELO DE TABELA -----------------------------
%Modelo da \autoref{tab:Instr3OP}:
%\begin{table}[htb]
%\centering
%\ABNTEXfontereduzida
%\caption{Exemplo de um formato de Instruções} \label{tab:Instr3OP}
%\begin{tabular}{r|p{2cm}|p{1.7cm}|p{1.7cm}|p{1.7cm}|p{3cm}} 
%\textbf{Tamanho (bits)} & 6 & 5 & 5 & 5 & 11 \\ \hline
%\textbf{Campo} & OPcode & Reg1 & Reg2 & Reg3 & Blank \\ \hline
%\textbf{Bits} & 31 - 26 & 25 - 21 & 20 - 16 & 15 - 11 & 10 - 0\\
%\end{tabular}
%\legend{Fonte: O Autor}
%\end{table}
%\end

%Como utilizar \textbf{Anexos}
%Material que não seja de autoria própria que traga maiores detalhes de algo utilizado no decorrer do relatório

%Como utilizar \textbf{Apêndices}
%Material de autoria própria que traga maiores detalhes e que não seja %essencial no decorrer do relatório.

%Modos de comandos $\backslash$ref\\
%$\backslash$ref: \ref{fig:esquematico}\\
%$\backslash$nameref: \nameref{fig:esquematico}\\
%$\backslash$autoref: \autoref{fig:esquematico}\\

%Existe uma forma de inserir arquivos de códigos diretamente no latex, verificar a formatação nas linhas 70 a 93 do código latex e a inserção feita do código \nameref{code:BinToBCD2.v} no \autoref{code:BinToBCD2.v} 

\chapter{Implementação}

\section{Entrega}
\par Abaixo podemos encontrar o link para o repositório contendo as implementações solicitadas: \url{https://github.com/parisi18/Calculo-Numerico}.

\section{Métodos}
\par A implementação dos métodos pode ser encontrado no arquivo metodos.py do repositório, que pode ser acessado via o link: \url{https://github.com/parisi18/Calculo-Numerico/blob/main/metodos.py}

\chapter{Exercícios didáticos}

\section{Isolamento}

\par Escreva tabelar\_sinais(f, a, b, n) que avalia f em n pontos igualmente espaçados e retorna
a lista de intervalos onde há mudança de sinal.

\par (a) Aplique a $f(x) = x^3 - 9x + 3$ em $[-5, 5]$ com $n=21$, $n=11$, $n=6$ e $n=4$. As quatro malhas encontram o mesmo número de raízes?

\begin{figure}[h]
    \centering
    \includegraphics[width=1\textwidth]{images/2.1a.png}
    \caption{Saída para exercício 2.1a}
    \label{fig:2.1a.png}
\end{figure}

\par (b) Agora aplique a $g(x) = (x - 1,05)(x - 1,15)(x - 3)$ em $[0, 4]$, com $n=9$, $n=17$, $n=41$ e $n=401$. Quantas raízes cada malha encontra? Compare com o número de raízes que você sabe que g tem.

\begin{figure}[h]
    \centering
    \includegraphics[width=1\textwidth]{images/2.1b.png}
    \caption{Saída para exercício 2.1b}
    \label{fig:2.1b.png}
\end{figure}

\par (c) No relatório: qual característica de g faz o tabelamento grosseiro falhar, e por que f não
sofre do mesmo problema? Que estratégia você adotaria se não soubesse de antemão onde estão
as raízes?

\begin{figure}[h]
    \centering
    \includegraphics[width=1\textwidth]{images/2.1c.png}
    \caption{Saída para exercício 2.1c}
    \label{fig:2.1c.png}
\end{figure}

\section{Previsão × realidade na bissecção}

\par Para cada $\varepsilon \in \{10^{-2}, 10^{-4}, 10^{-6}, 10^{-8}, 10^{-10}\}$, no intervalo $[0, 1]$:

\begin{itemize}
    \item Calcule o número previsto de iterações pela fórmula 
    \[
    k > \frac{\log(b_0 - a_0) - \log \varepsilon}{\log 2};
    \]
    \item Rode seu código e registre o número efetivo;
    \item Monte uma tabela comparando os dois.
\end{itemize}

\par A previsão é exata, otimista ou pessimista? Por quê?

\begin{figure}[h]
    \centering
    \includegraphics[width=1\textwidth]{images/2.2.png}
    \caption{Saída para exercício 2.2}
    \label{fig:2.2.png}
\end{figure}

\section{Custo real: avaliações de função}

Encontre a raiz em $(0, 1)$ com $\varepsilon = 10^{-8}$ pelos três métodos e preencha:

\begin{center}
\begin{tabular}{lccc}
\hline
\textbf{Método} & \textbf{Iterações} & \textbf{Avaliações de $f$} & \textbf{Avaliações de $f'$} \\ \hline
Bissecção & \rule{2cm}{0.4pt} & \rule{2cm}{0.4pt} & \rule{2cm}{0.4pt} \\
Newton    & \rule{2cm}{0.4pt} & \rule{2cm}{0.4pt} & \rule{2cm}{0.4pt} \\
Secante   & \rule{2cm}{0.4pt} & \rule{2cm}{0.4pt} & \rule{2cm}{0.4pt} \\ \hline
\end{tabular}
\end{center}

\par Discuta: \textbf{qual método é o mais barato em avaliações de função?} A resposta é a mesma que ``qual usa menos iterações''?

\begin{figure}[h]
    \centering
    \includegraphics[width=1\textwidth]{images/2.3.png}
    \caption{Saída para exercício 2.3}
    \label{fig:2.3.png}
\end{figure}

\section{Ordem empírica de convergência}

Sabendo que a raiz exata é $\xi = 0,3376089559658377$, calcule o erro $e_k = |x_k - \xi|$ a cada iteração e estime a ordem de convergência por
\[
p_k \approx \frac{\ln(e_{k+1}/e_k)}{\ln(e_k/e_{k-1})}
\]

\par Faça isso para Newton e para a secante. Compare com os valores teóricos $p = 2$ e $p \approx 1,618$.

\begin{figure}[h]
    \centering
    \includegraphics[width=1\textwidth]{images/2.4.png}
    \caption{Saída para exercício 2.4}
    \label{fig:2.4.png}
\end{figure}

\section{Os modos de falha de Newton}

Investigue computacionalmente cada situação abaixo. Para cada uma, \textbf{mostre a saída do seu código} e explique geometricamente o que aconteceu:

\begin{center}
\begin{tabular}{lccl}
\hline
\textbf{Caso} & \textbf{Função} & \textbf{$x_0$} & \textbf{O que investigar} \\ \hline
(a) & $f(x) = x^3 - 2x + 2$ & $0$ & Rode 10 iterações. Converge? \\
(b) & $f(x) = \arctan(x)$ & $2,0$ & E com $x_0 = 1,0$? Encontre o valor-limite de $x_0$ \\
(c) & $f(x) = x^3 - 9x + 3$ & $\sqrt{3}$ & O que sua implementação faz? \\ \hline
\end{tabular}
\end{center}

\begin{figure}[h]
    \centering
    \includegraphics[width=1\textwidth]{images/2.5a.png}
    \caption{Saída para exercício 2.5a}
    \label{fig:2.5a.png}
\end{figure}

\begin{figure}[h]
    \centering
    \includegraphics[width=1\textwidth]{images/2.5b.png}
    \caption{Saída para exercício 2.5b}
    \label{fig:2.5c.png}
\end{figure}

\begin{figure}[h]
    \centering
    \includegraphics[width=1\textwidth]{images/2.5c.png}
    \caption{Saída para exercício 2.5c}
    \label{fig:2.5c.png}
\end{figure}

\section{Raiz múltipla} 

Considere $f(x) = (x - 2)^2(x + 1)$, que tem raiz \textbf{dupla} em $x = 2$.

\begin{itemize}
    \item Rode Newton a partir de $x_0 = 3$ por 10 iterações e tabele o erro. A convergência é quadrática?
    \item Calcule a razão $e_{k+1}/e_k$. Ela se aproxima de que valor? Relacione com a multiplicidade $m = 2$.
    \item Implemente o \textbf{Newton modificado}, $x_{k+1} = x_k - m \cdot f(x_k)/f'(x_k)$, e refaça. O que mudou?
\end{itemize}

\begin{figure}[h]
    \centering
    \includegraphics[width=1\textwidth]{images/2.6.png}
    \caption{Saída para exercício 2.6}
    \label{fig:2.6.png}
\end{figure}

\section{A armadilha do resíduo}

Considere $f(x) = (x - 1)^{10}$.

\begin{itemize}
    \item Calcule $f(1,1)$ e $f(1,3)$. Comente a ordem de grandeza.
    
    \item Rode a bissecção em $[0; 1,5]$ com critério de parada \textbf{apenas} $|f(x)| < 10^{-8}$. Qual raiz você obtém? Qual o erro real em $x$?
    
    \item Refaça com o critério $|x_{k+1} - x_k| < 10^{-8}$. Compare.
    
    \item \textbf{Conclusão para o relatório:} em que situação o critério do resíduo é perigoso? E existe alguma situação em que ele é o critério \textit{mais} adequado?
\end{itemize}

\begin{figure}[h]
    \centering
    \includegraphics[width=1\textwidth]{images/2.7.png}
    \caption{Saída para exercício 2.7}
    \label{fig:2.7.png}
\end{figure}

\chapter{Situações-problema}

\begin{tcolorbox}[colback=gray!5!white, colframe=gray!20!white, arc=0mm, boxrule=0.5pt]
Em \textbf{todos} os problemas desta parte você deve, obrigatoriamente:
\begin{enumerate}
    \item Deduzir e escrever a função $f(x)$ cuja raiz resolve o problema.
    \item \textbf{Fazer a Fase I explicitamente} --- mostrar como isolou a raiz (tabelamento ou gráfico).
    \item Justificar a escolha do método e do chute inicial.
    \item Reportar o resultado \textbf{com unidade física} e com número de algarismos significativos coerente com os dados.
    \item Verificar a resposta substituindo de volta no problema original.
\end{enumerate}
\end{tcolorbox}

\section*{Problema A \textnormal{---} Reservatório esférico}

\par Um reservatório \textbf{esférico} de raio $R = 3,0\text{ m}$ é usado para armazenar água. O volume de líquido em função da altura $h$ medida a partir do fundo é
\[
V(h) = \pi h^2 \frac{3R - h}{3}
\]

\textbf{A.1} Determine a altura $h$ para que o reservatório contenha $40\text{ m}^3$ de água, com erro inferior a $1\text{ mm}$.

\textbf{A.2} A equação resultante é uma cúbica. Encontre \textbf{todas as três raízes reais} (use os três métodos, ou a bissecção em três intervalos distintos). Duas delas não servem. \textbf{Explique o que cada raiz espúria significa} e qual restrição física as elimina.

\textbf{A.3} Um sensor de nível será instalado. Construa uma tabela $h \times V$ para $V = 10, 20, \dots, 110\text{ m}^3$ resolvendo a equação para cada volume. Plote $h$ em função de $V$. \textit{Pergunta:} por que a curva é mais ``achatada'' no meio do que nas pontas?

\section*{3 - Situações-Problema}

\subsection*{A.1 Reservatório Esférico}

\textbf{Substituindo os valores na fórmula:}

\[
40 = \frac{\pi h^2(3(3)-h)}{3}
\]

\[
40 = \frac{\pi h^2(9-h)}{3}
\]

\[
120 = \pi h^2(9-h)
\]

\[
120 = 9\pi h^2-\pi h^3
\]

\textbf{Reorganizando a fórmula:}

\[
f(h)=\pi h^3-9\pi h^2+120=0
\]

Dividindo por $\pi$:

\[
g(h)=h^3-9h^2+\frac{120}{\pi}
\]

\[
g(h)=h^3-9h^2+38,1972
\]

\textbf{Critério de parada}

\[
|h_{k+1}-h_k|<0,001,
\qquad 0\leq h\leq 2R=6\,m
\]

Considerando o volume total do reservatório:

\[
V(6)=\frac{4}{3}\pi(3)^3=36\pi\approx113,05\,m^3
\]

Como $40\,m^3$ é aproximadamente $\frac{1}{3}$ do volume, deve estar próximo de $2$ e acima de $3$.

\[
g(2)=(2)^3-9(2)^2+38,1972
\]

\[
g(2)=10,1972>0
\]

\[
g(3)=(3)^3-9(3)^2+38,1972
\]

\[
g(3)=-15,8028<0
\]


\textbf{Chute inicial:}
Como a raiz está entre 2 e 3, o valor sera:

\[
h_0=2,5\,m
\]

\textbf{Método a ser utilizado:}
Neste exercício, o método de Newton-Raphson foi utilizado, pois já se tinha uma boa aproximação inicial, além de possuir uma $f'(x)$ fácil de calcular.

\begin{figure}[H]
    \centering
    \includegraphics[width=1\textwidth]{images/3A1_codigo.png}
    \caption{Saída Problema 3 A.1}
    \label{fig:3A1_codigo.png}
\end{figure}

\begin{figure}[H]
    \centering
    \includegraphics[width=1\textwidth]{images/3A2.png}
    \caption{Resposta Problema 3 A.2}
    \label{fig:3A2.png}
\end{figure}

\subsection*{A2}

Usando o Teorema de Bolzano (Teorema do Valor Intermediário): Se $g(h)$ muda de sinal nos extremos de um intervalo $[a,b]$, então existe pelo menos uma raiz real dentro desse intervalo.

\[
I_1=[-3,0], \qquad h_1=-1,87\,m
\]

\[
I_2=[2,3], \qquad h_2=2,407\,m
\quad\text{(número negativo)}
\]

\[
I_3=[6,5], \qquad h_3=6,467\,m
\quad\text{(Acima da esfera, já que }0\leq h\leq6,0\,m\text{)}
\]


\begin{figure}[H]
    \centering
    \includegraphics[width=1\textwidth]{images/3A3_codigo.png}
    \caption{Saída Problema 3 A.3}
    \label{fig:3A3_codigo.png}
\end{figure}

\section*{Problema B \textnormal{---} Perda de carga em tubulação}

O fator de atrito $f$ do escoamento turbulento em tubos é dado implicitamente pela \textbf{equação de Colebrook-White}:
\[
\frac{1}{\sqrt{f}} = -2 \log_{10}\left( \frac{\epsilon}{3,7 D} + \frac{2,51}{Re \sqrt{f}} \right)
\]

Dados de um trecho de adutora:
\begin{center}
\begin{tabular}{lr}
\hline
\textbf{Grandeza} & \textbf{Valor} \\ \hline
Diâmetro interno $D$ & $0,100\text{ m}$ \\
Rugosidade absoluta $\epsilon$ & $4,5 \times 10^{-5}\text{ m}$ (aço comercial) \\
Número de Reynolds $Re$ & $2,0 \times 10^5$ \\
Comprimento $L$ & $500\text{ m}$ \\
Vazão $Q$ & $0,050\text{ m}^3\text{/s}$ \\ \hline
\end{tabular}
\end{center}

\textbf{B.1} Determine $f$ com 6 casas decimais.

\par Para definirmos f, podemos assumir:
\[
F(x) = x + 2 \log_{10}\left( \frac{\epsilon}{3,7 D} + \frac{2,51x}{Re} \right)
\]
onde $x = \frac{1}{\sqrt{f}}$.

\par Assim, o problema passa a ser primeiramente achar x que zera aquela função. E, como a primeira derivada $F'(x)$ seria muito difícil de se achar para essa função, descartamos Newton-Raphson e a Secante possui ordem de convergência superior à da bissecção, utilizamos o método da Secante.

\begin{figure}[H]
    \centering
    \includegraphics[width=1\textwidth]{images/3B1.png}
    \caption{Saída Problema 3 B.1}
    \label{fig:3B1.png}
\end{figure}

\textbf{B.2} Este é o problema em que a \textbf{secante brilha}. Explique por quê: escreva (ou tente escrever) $f'$ analiticamente e comente a dificuldade. Resolva pelos três métodos e compare o esforço.

\par Como dito anteriormente ele possui uma boa ordem de convergência (superior a da bissecção) e 

\textbf{B.3} Um bom chute inicial vem da fórmula explícita de \textbf{Swamee-Jain}:
\[
f_0 = \frac{0,25}{\left[ \log_{10}\left( \frac{\epsilon}{3,7 D} + \frac{5,74}{Re^{0,9}} \right) \right]^2}
\]
\section*{Problema C \textnormal{---} Equação de van der Waals}

O gás ideal falha em altas pressões. A equação de van der Waals corrige isso:
\[
\left(P + \frac{a}{v^2}\right) (v - b) = RT
\]

onde $v$ é o volume molar ($\text{m}^3/\text{mol}$) e $R = 8,314\text{ J}/(\text{mol}\text{ K})$.

Para o $\text{CO}_2$: $a = 0,3640\text{ Pa}\text{ m}^6/\text{mol}^2$ e $b = 4,267 \times 10^{-5}\text{ m}^3/\text{mol}$.

\textbf{C.1} Para $T = 300\text{ K}$ e $P = 5,0\text{ MPa}$, calcule $v$ pelo modelo de gás ideal ($v = RT/P$) e depois pela equação de van der Waals. Qual o \textbf{erro percentual} do modelo ideal?

\textbf{C.2} Rearranjando, obtém-se a forma polinomial
\[
P v^3 - (Pb + RT) v^2 + a\,v - ab = 0
\]

Investigue quantas raízes \textbf{reais} existem nas condições do item C.1. Use tabelamento em uma faixa ampla ($10^{-5}$ a $1 \times 10^{-2}\text{ m}^3/\text{mol}$) --- atenção, a escala é muito pequena, considere usar espaçamento logarítmico.

\textbf{C.3} \textit{(discussão)} Abaixo da temperatura crítica ($T_c = 304,2\text{ K}$ para o $\text{CO}_2$) e em certas pressões, o cúbico tem \textbf{três} raízes reais positivas. Pesquise e explique brevemente o significado físico de cada uma. Qual delas seu método encontraria e por quê?

\textbf{C.4} Construa a \textbf{isoterma} $P \times v$ para $T = 300\text{ K}$, variando $P$ de $1$ a $10\text{ MPa}$ em passos de $0,5\text{ MPa}$. Plote e compare com a curva do gás ideal no mesmo gráfico.

{Passo 1 --- Dedução de $f(v)$.}
Queremos o volume molar $v$ que satisfaz a equação de van der Waals para $T=300$ K e $P=5{,}0$ MPa. Multiplicando $\left(P + \dfrac{a}{v^2}\right)(v-b) = RT$ por $v^2$ e rearranjando, obtemos a forma polinomial já apresentada no enunciado (C.2), cuja raiz é o volume molar procurado:
\[
f(v) = P v^3 - (Pb + RT)v^2 + a\,v - ab = 0 .
\]
Como referência de comparação (modelo de gás ideal), calculamos também
\[
v_{\text{ideal}} = \frac{RT}{P} = \frac{8{,}314 \times 300}{5{,}0\times10^{6}} = 4{,}988\times10^{-4}\ \text{m}^3/\text{mol}.
\]

{Passo 2 --- Fase I: isolamento da raiz.}
Como os valores de $v$ são da ordem de $10^{-4}\ \text{m}^3/\text{mol}$, um tabelamento linear seria grosseiro demais ou exigiria um número impraticável de pontos. Por isso, tabelamos $f(v)$ com \textbf{espaçamento logarítmico} em $[10^{-5},\,10^{-2}]\ \text{m}^3/\text{mol}$ (400 pontos), conforme pedido no item C.2. O tabelamento encontrou uma \textbf{única} mudança de sinal nessa faixa, isolando a raiz no intervalo
\[
v \in [\,3{,}601\times10^{-4},\ 3{,}664\times10^{-4}\,]\ \text{m}^3/\text{mol},
\]
próximo de $v_{\text{ideal}}$, como esperado fisicamente (fase vapor).

{Passo 3 --- Justificativa do método e do chute inicial.}
Escolhemos o método de \textbf{Newton-Raphson}. A derivada
\[
f'(v) = 3Pv^2 - 2(Pb+RT)v + a
\]
é um polinômio simples e barato de obter analiticamente (sem risco de erro de derivação), e Newton converge quadraticamente quando o chute inicial é bom. Usamos $x_0 = v_{\text{ideal}} = 4{,}988\times10^{-4}\ \text{m}^3/\text{mol}$, pois na faixa de pressão considerada o comportamento real do gás não deve se afastar muito do modelo ideal, e o Passo 2 confirmou que a raiz física está de fato próxima desse valor.

{Passo 4 --- Resultado.}
Com critério de parada $|f(v_{k+1})|<\varepsilon$ ou $|v_{k+1}-v_k|<\varepsilon$ (tolerância relativa, $\varepsilon = 10^{-6}\,v_{\text{ideal}}$), o método convergiu em $5$ iterações para
\[
v_{\text{vdW}} = 3{,}656\times10^{-4}\ \text{m}^3/\text{mol} \quad (3\text{ algarismos significativos, coerentes com } a \text{ e } b).
\]
O erro percentual do modelo de gás ideal em relação a esse resultado é
\[
\frac{|v_{\text{ideal}} - v_{\text{vdW}}|}{v_{\text{vdW}}}\times100\% \approx 36{,}4\%,
\]
uma diferença expressiva, o que confirma que o modelo ideal falha nessas condições de alta pressão, como antecipado no enunciado.

{Passo 5 --- Verificação.}
Substituindo $v_{\text{vdW}}$ de volta na equação original:
\[
\left(P+\frac{a}{v_{\text{vdW}}^2}\right)(v_{\text{vdW}}-b) = 2494{,}20\ \text{J/mol} \approx RT = 2494{,}20\ \text{J/mol}.
\]
O resíduo do polinômio é $f(v_{\text{vdW}}) \approx 1{,}3\times10^{-13}$ ($\approx 0$), confirmando a raiz.

{C.2 --- Número de raízes reais.}
Como mostrado no Passo 2, nas condições de $T=300$ K e $P=5{,}0$ MPa existe apenas \textbf{uma} raiz real positiva no intervalo investigado $[10^{-5}, 10^{-2}]\ \text{m}^3/\text{mol}$, ou seja, $v_{\text{vdW}}=3{,}656\times10^{-4}\ \text{m}^3/\text{mol}$ é a única solução física relevante nessas condições.

{C.3 --- Discussão física das três raízes.}
A pressão de saturação do $\text{CO}_2$ a $300$ K é da ordem de $\sim 6{,}7$ MPa. Como $P = 5{,}0$ MPa está \textbf{abaixo} dessa saturação, o $\text{CO}_2$ está em fase \textbf{vapor pura} (região de uma só fase), o que explica a raiz única encontrada nos itens C.1/C.2. De forma geral, para $T<T_c=304{,}2$ K e $P$ \textbf{acima} da pressão de saturação, o cúbico de van der Waals pode ter até três raízes reais positivas (região de coexistência líquido--vapor da isoterma):
\begin{itemize}
    \item a \textbf{menor} raiz corresponde ao volume molar do \textbf{líquido saturado};
    \item a \textbf{maior} raiz corresponde ao volume molar do \textbf{vapor saturado};
    \item a raiz \textbf{intermediária} não tem significado físico (região onde $dP/dv>0$, mecanicamente instável --- um artefato da curva suave de van der Waals, que na isoterma real tem um patamar horizontal em vez dessa "corcova").
\end{itemize}
Partindo de um chute próximo do gás ideal (fase vapor), o método de Newton sempre converge para a raiz \textbf{maior} (fase vapor) quando ela existe, que é a de maior interesse prático ao tratar o $\text{CO}_2$ como fluido de processo em fase gasosa.

{C.4 --- Isoterma $P\times v$.}
A Figura~\ref{fig:isoterma_C} compara a isoterma de van der Waals com a do gás ideal para $T=300$ K, $P$ variando de $1$ a $10$ MPa em passos de $0{,}5$ MPa (cada ponto resolvido por Newton, com chute $v_0 = RT/P$ específico para cada pressão). Em toda a faixa, o gás ideal superestima o volume molar em relação ao van der Waals --- o termo de atração molecular $a/v^2$ "puxa" as moléculas para mais perto, reduzindo o volume real ocupado frente ao previsto pelo modelo ideal.

\begin{figure}[h!]
    \centering
    \includegraphics[width=0.7\textwidth]{images/problema_C_isoterma.png}
    \caption{Isoterma $P\times v$ em $T=300$ K: van der Waals \textit{versus} gás ideal.}
    \label{fig:isoterma_C}
\end{figure}

\section*{Problema D \textnormal{---} Taxa Interna de Retorno}

Uma empresa avalia um projeto com o seguinte fluxo de caixa (em mil R\$):

\begin{center}
\begin{tabular}{lccccc}
\hline
\textbf{Ano} & $0$ & $1$ & $2$ & $3$ & $4$ \\ \hline
\textbf{Fluxo} & $-1000$ & $+300$ & $+350$ & $+400$ & $+450$ \\ \hline
\end{tabular}
\end{center}

O \textbf{Valor Presente Líquido} a uma taxa $i$ é
\[
VPL(i) = \sum_{k=0}^{n} \frac{C_k}{(1 + i)^k}
\]

A \textbf{Taxa Interna de Retorno (TIR)} é a taxa que zera o VPL.

\textbf{D.1} Calcule a TIR com precisão de $10^{-6}$.

\textbf{D.2} Plote $VPL(i)$ para $i \in [0; 0,5]$ e mostre graficamente onde está a raiz.

\textbf{D.3} \textit{(decisão)} Se o custo de capital da empresa for $15\%$ ao ano, o projeto deve ser aceito? E se for $20\%$? Justifique com o VPL calculado em cada caso.

\textbf{D.4} \textit{(o problema das múltiplas TIRs)} Agora considere um segundo projeto:

\begin{center}
\begin{tabular}{lccc}
\hline
\textbf{Ano} & $0$ & $1$ & $2$ \\ \hline
\textbf{Fluxo} & $-1000$ & $+2500$ & $-1540$ \\ \hline
\end{tabular}
\end{center}

Plote o VPL e mostre que existem \textbf{duas} taxas que o zeram. Encontre ambas.

\begin{tcolorbox}[colback=red!5!white, colframe=red!20!white, arc=0mm, boxrule=0.5pt]
\textbf{Pergunta crítica.} O que acontece se um analista rodar Newton com um único chute inicial e reportar o resultado sem fazer a Fase I? Que decisão errada isso pode gerar?
\end{tcolorbox}

{Passo 1 --- Dedução de $f(i)$.}
A TIR é, por definição, a taxa $i$ que zera o VPL. Assim, a raiz procurada é a de
\[
f(i) = VPL(i) = \sum_{k=0}^{4} \frac{C_k}{(1+i)^k} = -1000 + \frac{300}{1+i} + \frac{350}{(1+i)^2} + \frac{400}{(1+i)^3} + \frac{450}{(1+i)^4} = 0.
\]
Para uso em Newton, precisamos também de
\[
f'(i) = \frac{d(VPL)}{di} = \sum_{k=0}^{4} \frac{-k\,C_k}{(1+i)^{k+1}}.
\]

{Passo 2 --- Fase I: isolamento da raiz.}
Tabelamos $f(i)=VPL(i)$ em $i\in[0;\,0{,}5]$ com $51$ pontos igualmente espaçados (item D.2). O tabelamento encontrou uma única mudança de sinal, isolando a raiz no intervalo
\[
i \in [0{,}17,\ 0{,}18],
\]
como mostra a Figura~\ref{fig:vpl_projeto1}: a curva $VPL(i)$ é monotonicamente decrescente em todo o intervalo $[0;0{,}5]$ (fluxo de caixa convencional --- uma única troca de sinal em $C_k$), cruzando o eixo uma única vez.

\begin{figure}[ht]
    \centering
    \includegraphics[width=0.7\textwidth]{images/problema_D_vpl_projeto1.png}
    \caption{$VPL(i)$ do projeto 1 em $i\in[0;50\%]$, com a raiz (TIR) destacada.}
    \label{fig:vpl_projeto1}
\end{figure}

{Passo 3 --- Justificativa do método e do chute inicial.}
Escolhemos o método de \textbf{Newton-Raphson}. Justificativa: $f'(i)$ é uma soma simples de potências de $(1+i)$, barata de calcular analiticamente, e a curva $VPL(i)$ é suave e sem inflexões bruscas na região isolada --- condições ideais para a convergência quadrática de Newton. Usamos o chute inicial $i_0 = 0{,}10$ ($10\%$), uma taxa de referência de mercado tipicamente usada como primeira estimativa e razoavelmente próxima da região já isolada no Passo 2.

{Passo 4 --- Resultado.}
Com critério de parada $|VPL(i_{k+1})|<\varepsilon$ ou $|i_{k+1}-i_k|<\varepsilon$, $\varepsilon=10^{-6}$ (conforme exigido em D.1), o método convergiu em $4$ iterações para
\[
\text{TIR} = 17{,}0937\%\ \big(i \approx 0{,}170937\big),
\]
reportada com 6 algarismos significativos, coerente com a tolerância $10^{-6}$ pedida.

{Passo 5 --- Verificação.}
Substituindo a TIR de volta na definição do VPL:
\[
VPL(0{,}170937) \approx 1{,}7\times10^{-12}\ \text{mil R\$} \approx 0,
\]
confirmando que a taxa encontrada de fato zera o valor presente líquido do fluxo de caixa original.

{D.3 --- Decisão de aceitar ou rejeitar o projeto.}
Calculando o VPL diretamente nos custos de capital pedidos:
\begin{itemize}
    \item Custo de capital $=15\%$: $VPL(0{,}15) = 45{,}82$ mil R\$ $>0$ $\Rightarrow$ \textbf{ACEITAR} o projeto (TIR $=17{,}09\%>15\%$).
    \item Custo de capital $=20\%$: $VPL(0{,}20) = -58{,}45$ mil R\$ $<0$ $\Rightarrow$ \textbf{REJEITAR} o projeto (TIR $=17{,}09\%<20\%$).
\end{itemize}
Em ambos os casos, a decisão é consistente com a regra clássica da TIR: aceitar se $\text{TIR} > $ custo de capital.

{D.4 --- O problema das múltiplas TIRs.}
Para o segundo projeto, fluxo $[-1000,\,+2500,\,-1540]$, há \textbf{duas} trocas de sinal no fluxo de caixa (não-convencional). Tabelando $VPL(i)$ em uma faixa ampla ($i\in[-0{,}3;\,1{,}0]$), identificamos \textbf{dois} intervalos com mudança de sinal, isolando as duas raízes em $[9{,}5\%;\,10{,}2\%]$ e $[39{,}6\%;\,40{,}2\%]$. Aplicando bissecção separadamente em cada intervalo (garantindo achar ambas as raízes, sem depender de um único chute), obtemos:
\[
\text{TIR}_1 = 10{,}0000\%, \qquad \text{TIR}_2 = 40{,}0000\%,
\]
ambas verificadas por $VPL(\text{TIR}_1)\approx -1{,}1\times10^{-6}$ e $VPL(\text{TIR}_2)\approx -3{,}3\times10^{-7}$ (mil R\$), praticamente nulos. A Figura~\ref{fig:vpl_projeto2} mostra a curva $VPL(i)$ do projeto 2, com as duas raízes destacadas.

\begin{figure}[ht]
    \centering
    \includegraphics[width=0.7\textwidth]{images/problema_D_vpl_projeto2.png}
    \caption{$VPL(i)$ do projeto 2, mostrando as duas TIRs (10\% e 40\%).}
    \label{fig:vpl_projeto2}
\end{figure}

{Resposta à pergunta crítica.}
Se um analista rodar Newton com um único chute inicial (por exemplo, $i_0=0{,}5$) e reportar o resultado sem fazer a Fase I, ele encontra \textbf{apenas uma} das duas TIRs (10\% ou 40\%, dependendo do chute) e pode concluir erroneamente que o projeto tem uma única taxa de retorno bem definida. Isso é perigoso porque a regra da TIR (``aceite se $\text{TIR}>$ custo de capital'') só faz sentido quando há uma única raiz: com fluxo de caixa não-convencional como este, pode haver várias taxas que zeram o VPL, e a decisão de aceitar/rejeitar passa a depender de \textbf{qual} TIR foi encontrada. Um analista que pula a Fase I pode reportar ``TIR$=40\%$, superior ao custo de capital de $15\%$, ACEITAR'' sem perceber que a $10\%$ o projeto também zera e que, para custos de capital \textbf{entre} as duas raízes, o VPL pode ser negativo --- levando a uma decisão de investimento errada.

\section*{Problema E \textnormal{---} Equação de Kepler \textnormal{(8 pontos)}}

A posição de um corpo em órbita elíptica é obtida resolvendo a \textbf{equação de Kepler}:
\[
M = E - e \sin E
\]

onde $M$ é a anomalia média (conhecida, proporcional ao tempo), $e$ é a excentricidade da órbita e $E$ é a anomalia excêntrica (a incógnita).

\textbf{E.1} Para a órbita do cometa Halley, $e = 0,967$. Calcule $E$ para $M = 0,2\text{ rad}$.

\textbf{E.2} \textit{(robustez)} Resolva para os três casos abaixo, usando Newton com chute inicial $E_0 = M$:
\begin{center}
\begin{tabular}{lcc}
\hline
\textbf{Caso} & $e$ & \textbf{$M$ (rad)} \\ \hline
(i)   & $0,10$ & $0,5$ \\
(ii)  & $0,90$ & $0,1$ \\
(iii) & $0,99$ & $0,01$ \\ \hline
\end{tabular}
\end{center}

Tabele o número de iterações de cada caso. \textbf{O que acontece com o esforço à medida que $e \to 1$?} Explique olhando para $f'(E) = 1 - e \cos E$ perto de $E = 0$.

\textbf{E.3} Um chute inicial melhor para órbitas muito excêntricas é $E_0 = M + e \sin M$. Refaça o caso (iii) com ele e compare.

\textbf{E.4} \textit{(comparação)} Para o caso (iii), rode também a bissecção em $[0; \pi]$. Ela converge? Em quantas iterações? Qual método você recomendaria para um software de rastreamento de satélites que precisa resolver essa equação \textbf{milhões de vezes por segundo}, e por quê?

\chapter{Relatório e discussão}

Responda de forma dissertativa, com base nos seus próprios resultados.

\textbf{4.1} Construa uma tabela-resumo com todos os problemas da Parte 3, indicando para cada um: o método escolhido, o motivo da escolha e o número de iterações gasto.

\textbf{4.2} Em qual problema da Parte 3 a \textbf{Fase I} foi mais decisiva? Descreva o que teria dado errado se você a tivesse pulado.

\textbf{4.3} Em qual problema a \textbf{secante} foi claramente superior a Newton? E existiu algum em que Newton foi claramente superior à secante? Fundamente com os números de avaliações de função que você mediu.

\textbf{4.4} Um colega afirma: \textit{``a bissecção é obsoleta, ninguém usa isso na prática''}. Escreva uma resposta técnica de até 10 linhas, usando pelo menos um resultado do seu próprio trabalho como argumento.

\textbf{4.5} Descreva, em pseudocódigo, como você combinaria bissecção e Newton em um \textbf{método híbrido} que seja simultaneamente seguro e rápido. Quando seu algoritmo decidiria trocar de um para o outro?

% ---
% Finaliza a parte no bookmark do PDF
% para que se inicie o bookmark na raiz
% e adiciona espaço de parte no Sumário
% ---
\phantompart


% ----------------------------------------------------------
% ELEMENTOS PÓS-TEXTUAIS
% ----------------------------------------------------------
\postextual

% ----------------------------------------------------------
% Referências bibliográficas
% ----------------------------------------------------------


% ----------------------------------------------------------
% Apêndices
% ----------------------------------------------------------

% ---
% Inicia os apêndices
% ---
%\begin{apendicesenv}

% Imprime uma página indicando o início dos apêndices
%\partapendices

% ----------------------------------------------------------
%\chapter{BinToBCD2.v}
%\label{code:BinToBCD2.v}
%Código do arquivo BinToBCD2.v
%\lstinputlisting{Codigos/BinToBCD2.v} 
% ----------------------------------------------------------

% ----------------------------------------------------------
%\chapter{Apêndice 2}
% ----------------------------------------------------------


%\end{apendicesenv}
% ---


% ----------------------------------------------------------
% Anexos
% ----------------------------------------------------------

% ---
% Inicia os anexos
% ---
%\begin{anexosenv}

% Imprime uma página indicando o início dos anexos
%\partanexos

% ---
%\chapter{Anexo 1}
% ---

% ---
%\chapter{Anexo 2}
% ---

%\end{anexosenv}

%---------------------------------------------------------------------
% INDICE REMISSIVO
%---------------------------------------------------------------------

\phantompart

\printindex


\end{document}
